\section{Methodology}
\label{sec:methodology}

\subsection{Research Design}

This work adopts a controlled experimental approach to compare three time-series databases
across nine distinct IoT workload patterns and three data-volume scales. The evaluation covers
write throughput, per-operation latency (average and 99th-percentile), and resource consumption
(CPU and memory) for each combination of database, workload type, and scale. The goal is to
characterize which architectural trade-offs favour which application profiles, going beyond
aggregate rankings to explain \textit{why} each database performs as it does.

A central design objective beyond performance measurement is \textit{reproducibility}.
The comparative literature reviewed in this work reveals a systemic gap: of the five recent
TSDB performance studies analyzed, only one published sufficient scripts and configurations
for independent replication~\citep{pereira2023mulletbench}. The others either omitted
configurations entirely, described their setups in prose too vague to follow mechanically,
or used toolchains that require institutional access~\citep{shah2022performance,naqvi2017influxdb,
lima2024avaliando,wang2023iotdb}. This work treats reproducibility as a first-class output:
every component needed to replicate the results is version-pinned and publicly available
alongside this paper.

\subsection{Database Selection}

The three databases were selected to represent distinct architectural approaches to
time-series storage, each identified as a leading system in the reviewed literature.

\textbf{InfluxDB 2.7} is the de facto reference standard in TSDB comparative research,
appearing in four of the five performance studies reviewed~\citep{naqvi2017influxdb,
shah2022performance,lima2024avaliando,wang2023iotdb}. It uses a TSM (Time-Structured Merge
Tree) storage engine that organizes data into compressed time-ordered blocks. Ingestion
occurs via an HTTP line-protocol API; analytics queries use the Flux expression language.
InfluxDB's ubiquity in prior benchmarks makes it an essential comparison anchor: any new
study that omits it loses comparability with the existing literature.

\textbf{Apache IoTDB 1.3.0}~\cite{wang2023iotdb} emerged as the performance leader in
the two most recent comparative studies~\citep{pereira2023mulletbench,wang2023iotdb},
representing a shift from the InfluxDB dominance reported in older benchmarks.
Developed at Tsinghua University specifically for industrial IoT, it uses a custom columnar
file format (TsFile) and a binary Thrift RPC protocol (SESSION\_BY\_TABLET). Write ingestion
is buffered in a JVM-managed heap before being flushed to disk as sorted columnar segments,
enabling exceptionally high throughput for ordered time-series streams. Its rapid rise in
recent benchmarks, combined with its distinct architectural approach, made it a natural
candidate for this study.

\textbf{TimescaleDB} extends PostgreSQL~15 with time-series capabilities: automatic
hypertable partitioning by time interval, the \texttt{time\_bucket()} aggregation function
for downsampling, and optional chunk-level compression. It retains the full relational engine
with WAL-based durability, MVCC concurrency control, B-tree indexes, and complete SQL
compatibility. For organisations with existing PostgreSQL infrastructure, or for applications
requiring joins between time-series and relational data, it provides a path to time-series
capability without abandoning the relational ecosystem. It appears in four of the six
reviewed benchmarks~\citep{shah2022performance,lima2024avaliando,wang2023iotdb,grzesik2020comparative}.

Plain PostgreSQL without the TimescaleDB extension was explicitly excluded. TimescaleDB
tends to outperform plain PostgreSQL in this domain; including both would duplicate tests.
Conclusions about general PostgreSQL configuration, such as \texttt{shared\_buffers}
sizing, are transferable, but results that depend on TimescaleDB-specific optimizations,
such as hypertables and chunk-level min/max pruning, do not apply directly to plain
PostgreSQL without the extension.

\subsection{Benchmark Tool}
\label{sec:benchmark-tool}

The workload generator is \textit{iot-benchmark}~\cite{iotbenchmark}, developed at Tsinghua
University, which ships purpose-built benchmark clients for all three databases under evaluation.
This native multi-database support means each workload is expressed in the correct protocol
for each target: the Thrift SESSION\_BY\_TABLET interface for IoTDB, the HTTP line-protocol
for InfluxDB, and a JDBC connection for TimescaleDB. A consistent, parameterized interface
ensures that differences in measured performance reflect the databases themselves rather
than the benchmarking layer.

Alternative tools were evaluated. The Time Series Benchmark Suite (TSBS) offers a mature
multi-database framework but lacks native IoTDB support at the time of this writing.
Apache JMeter, used in Lima et al.~\cite{lima2024avaliando}, handles HTTP-based systems well
but requires custom scripting for Thrift and JDBC interfaces, introducing potential tool-specific
overhead. iot-benchmark's uniform interface and IoT-specific workload semantics (concurrent clients
  organized into device and sensor hierarchies, configurable batch sizes, and
Poisson-distributed out-of-order data generation) made it the most appropriate choice
for this study.

The tool produces two structured output matrices that are parsed by the orchestration layer.
The Result Matrix reports total operations, elapsed time, and aggregate throughput (pts/s).
The Latency Matrix reports average, minimum, maximum, and percentile latencies per operation
type, providing the per-operation tail-behavior data needed to evaluate real-time IoT workloads.

\subsection{Workload Design}
\label{sec:workload}

The benchmark covers nine workload types across two categories, designed to span the range
of patterns encountered in production IoT systems.

\subsubsection{Write Tests}

Four write tests exercise distinct ingestion patterns:

\begin{description}
  \item[\texttt{write}] Baseline sequential ingestion with ordered timestamps and the
    default batch size of 100 rows per request. This measures raw sustained write throughput
    under ideal, ordered conditions.

  \item[\texttt{out-of-order}] Twenty percent of data points arrive with timestamps older
    than the current ingestion frontier, modeled as a Poisson-distributed delay. This simulates
    IoT sensors that lose connectivity and later replay buffered readings, a ubiquitous
    pattern in industrial and field deployments where sensors operate over unreliable cellular
    or LPWAN networks.

  \item[\texttt{batch-small}] \texttt{BATCH\_SIZE=1}: one data point per write request.
    This isolates per-request protocol overhead: an HTTP round-trip plus line-protocol parsing
    for InfluxDB, a full PostgreSQL client-server cycle plus WAL flush for TimescaleDB, and a
    single Thrift RPC call for IoTDB.

  \item[\texttt{batch-large}] \texttt{BATCH\_SIZE=1000}: 1{,}000 data points per request.
    This tests bulk-loading efficiency and reveals the practical ceiling of each database's
    batching architecture. The ratio between \texttt{batch-small} and \texttt{batch-large}
    throughput directly quantifies how much per-request overhead dominates over data-transfer cost.
\end{description}

Each write test clears and re-seeds the database before running to ensure independent,
history-free measurements.

\subsubsection{Read Tests}

Five read tests cover the query patterns most common in IoT analytics pipelines:

\begin{description}
  \item[\texttt{latest-point}] Retrieves the most recent value for each sensor. This is
    the foundational real-time monitoring query: a dashboard or alerting system continuously
    polling current device state.

  \item[\texttt{downsample}] Executes \texttt{GROUP~BY} time-bucket queries computing
    averages over fixed time windows. This models the standard dashboard aggregation use case,
    where raw sensor readings are collapsed into minute- or hour-level summaries.

  \item[\texttt{range-query}] Retrieves all data points within a specified time window,
    without filtering on values. This models historical data export, replay pipelines, and
    model retraining ingestion.

  \item[\texttt{value-filter}] Applies value predicates such as \texttt{temperature~>~85.0},
    including aggregations over value-filtered ranges. This models threshold alerting:
    ``find all sensor readings exceeding the alarm level in the past hour.'' Value-filter is
    the most architecturally discriminating query type in this benchmark because it cannot be
    answered from time-based indexes alone: each database must either maintain a secondary
    index or scan all stored values.

  \item[\texttt{read}] Runs all ten query types supported by iot-benchmark with equal weight,
    providing a comprehensive mixed-workload baseline.
\end{description}

All read tests operate on the dataset loaded by the preceding baseline \texttt{write} test.

\subsection{Scale Parameters}
\label{sec:scale-params}

Three data-volume scales are defined to expose how performance characteristics evolve with
dataset size. Table~\ref{tab:scales} summarizes the parameters.

\begin{table*}[t]
  \centering
  \caption{Benchmark scale parameters. All scales use a default batch size of 100 rows per
  request. Large scale contains approximately one hundred times more data than small.}
  \label{tab:scales}
  \begin{tabular}{lrrrrr}
    \hline
    \textbf{Scale} & \textbf{Clients} & \textbf{Devices} & \textbf{Sensors} & \textbf{Loops} & \textbf{Duration} \\
    \hline
    Small  &  5 & 10 & 10 & 1{,}000 & ${\sim}$5 min \\
    Medium &  5 & 10 & 10 & 5{,}000 & ${\sim}$30 min \\
    Large  & 10 & 50 & 20 & 5{,}000 & ${\sim}$2 h \\
    \hline
  \end{tabular}
\end{table*}

Large scale contains approximately 20$\times$ more data than medium and one hundred times more
than small, driven by both a higher loop count and an expanded topology (10 clients
$\times$ 50 devices $\times$ 20 sensors versus 5~$\times$~10~$\times$~10). Multi-scale
measurement is essential because several key characteristics (JVM JIT compilation
saturation, buffer exhaustion, checkpoint stall patterns, and index scan scaling) only
manifest or become dominant at higher data volumes. A single-scale benchmark risks drawing
conclusions that apply only to a narrow data range.

\subsection{Infrastructure and Reproducibility}
\label{sec:infrastructure}

Reproducibility in this project rests on three components.

\textbf{Docker Compose}~\cite{docker} manages the database containers with pinned image
versions: \texttt{apache/iotdb:1.3.0-standalone}, \texttt{influxdb:2.7-alpine}, and
\texttt{timescale/timescaledb:latest-pg15}. Containers are started and stopped
\textit{sequentially}: only one database is active at any given time. This design, introduced
after the initial runs revealed cross-database resource contention artefacts, gives each
database exclusive access to the host's CPU, memory bandwidth, and I/O.

\textbf{Nix flake}~\cite{nixos} is an optional tool that provides a declarative,
lock-file-pinned development environment. Running \texttt{nix~develop} drops into a shell
with exactly-specified versions of Java~17, Maven, Python~3, and Docker tooling. Running
\texttt{nix~run~.\#build} compiles all three benchmark clients in a single step. Every
transitive build dependency is recorded in the flake lock file, eliminating the
environment-variation failures that typically plague benchmark replications. Users without
Nix can replicate the experiments by manually installing the listed dependencies (Java~17,
Maven, Python~3, and Docker) on any compatible Linux host.

\textbf{Python CLI} (\texttt{benchmark.py}) orchestrates the complete benchmark workflow
programmatically. It starts the target container, writes the correct parameters to
the benchmark client's \texttt{conf/config.properties}, invokes the client, polls
\texttt{docker~stats} at one-second intervals to collect CPU and memory samples, and appends
a structured result row to the output table after each test. A complete nine-test,
three-database run at any scale reduces to:

\begin{verbatim}
  python3 benchmark.py --scale large
\end{verbatim}

This level of automation makes the results traceable to a specific, versioned execution
environment. Any researcher with the repository and Docker available can reproduce the full
dataset from scratch.

\subsection{Replication Package}
\label{sec:replication}

The complete benchmark infrastructure is publicly available at
\url{https://github.com/LuizFerK/iot-benchrunner}. Researchers can
reproduce the full experimental environment on any Linux host using
either the Nix toolchain or standard system packages.

\textbf{Requirements.} Java~17, Maven, Python~3, and Docker with Docker
Compose are required. Nix users obtain all dependencies automatically via
the provided flake; non-Nix users should install the same versions
manually.

\textbf{Setup.} Clone the repository including its \texttt{iot-benchmark}
submodule, then build the Java benchmark clients once:

\begin{verbatim}
git clone --recurse-submodules \
  https://github.com/LuizFerK/iot-benchrunner
cd iot-benchrunner

# Nix (recommended)
nix run .#build

# Non-Nix
cd iot-benchmark
mvn -pl influxdb-2.0,timescaledb,iotdb-1.3 \
  package -DskipTests -q
cd ..
\end{verbatim}

\textbf{Running.} A complete nine-test, three-database run at any scale
is a single command. Nix users should first enter the dev shell with
\texttt{nix~develop}, which provides all runtime dependencies.

\begin{verbatim}
python3 benchmark.py --scale small   # ~5 min
python3 benchmark.py --scale medium  # ~30 min
python3 benchmark.py --scale large   # ~2 h
\end{verbatim}

\textbf{Generating charts.} After the benchmark produces result CSVs in
the \texttt{results/} directory, the visualization script regenerates all
figures used in this paper:

\begin{verbatim}
python3 charts.py --source all \
  --language en-us
\end{verbatim}

The \texttt{--source} flag accepts individual scales (\texttt{small},
\texttt{medium}, \texttt{large}), rerun variants, and \texttt{comparison}
for cross-scale overlay plots. Output PNG files are written to the
\texttt{charts/} directory.

\subsection{Metrics}
\label{sec:metrics}

Table~\ref{tab:metrics} describes the performance metrics collected for each run.

\begin{table*}[t]
  \centering
  \caption{Performance metrics collected per benchmark run.}
  \label{tab:metrics}
  \begin{tabular}{ll}
    \hline
    \textbf{Metric} & \textbf{Source} \\
    \hline
    Throughput (pts/s)    & iot-benchmark Result Matrix \\
    Average latency (ms)  & iot-benchmark Latency Matrix \\
    P99 latency (ms)      & iot-benchmark Latency Matrix \\
    Avg/Peak CPU (\%)     & \texttt{docker stats} (1\,s interval) \\
    Avg/Peak Memory (MB)  & \texttt{docker stats} RSS \\
    Wall-clock time (s)   & Run duration \\
    \hline
  \end{tabular}
\end{table*}

For write tests, throughput measures data points ingested per second. For read tests,
throughput sums point counts returned across all active query types. Because different query
types return fundamentally different amounts of data per call: a range query returns
hundreds of points, a latest-point query returns exactly one. \textit{Latency} is therefore
used as the primary read comparison metric throughout Section~\ref{sec:experiments}.

P99 latency captures tail behavior: 99\% of operations completed faster than this value.
The gap between average and P99 reveals response-time consistency. A large gap signals
occasional severe outliers that would produce visible jitter in live IoT monitoring applications,
even when average latency appears acceptable.

\subsection{Configuration Tuning}
\label{sec:tuning}

After completing the initial benchmark runs, two methodological improvements were applied
before re-running all scales:

\begin{enumerate}
  \item \textbf{Sequential container lifecycle.} The original setup started all three
    database containers simultaneously. Although each test ran on one database at a time, idle
    containers competed for OS page cache and I/O bandwidth~\citep{shapira2016pitfalls}. TimescaleDB accumulated several
    gigabytes of page cache during IoTDB's write sessions, producing deeply misleading memory
    measurements. IoTDB's JVM pre-allocated its maximum heap during idle time for the same reason.
    The corrected setup starts each container immediately before its tests and stops it
    immediately after.

  \item \textbf{Database-level tuning.} TimescaleDB received a pgtune~\cite{pgtune}
    configuration for PostgreSQL~15 targeting this machine's hardware:
    \texttt{shared\_buffers=8GB}, \texttt{work\_mem=27413kB},
    \texttt{maintenance\_work\_mem=2GB}, \texttt{effective\_io\_concurrency=1000}, and
    \texttt{checkpoint\_completion\_target=0.9}. These settings are the standard baseline for
    any production PostgreSQL deployment on this class of hardware; the defaults are tuned
    for minimal-RAM environments and produce artificially poor performance on capable machines.
    InfluxDB's TSM write-cache ceiling was raised from the default 1\,GB to 8\,GB to match the
    available RAM. A 28\,GB memory limit was applied to every container.
\end{enumerate}

Both the original and tuned results are reported and analyzed at small scale, where the
contrast is most instructive. Medium and large-scale results use the tuned methodology
throughout. The comparison between original and tuned runs serves a dual purpose: it
quantifies the magnitude of contention artefacts and validates which performance
characteristics are genuine architectural properties of each database.
